\documentclass[11pt,a4paper]{article}

\usepackage[utf8]{inputenc}
\usepackage[french]{babel}
\usepackage{geometry}
\usepackage{amsmath,amssymb}
\usepackage{graphicx}
\usepackage{booktabs}
\usepackage{multirow}
\usepackage{hyperref}

\geometry{margin=2.5cm}

\title{%
    \textbf{Métaheuristiques pour le Set Packing Problem} \\
    \large Exercices d'Implémentation
}
\author{NAIT SLIMANI Maroua \& MOLLI Lila \& ABID Ikram \\ Master Informatique \\ Nantes Université}
\date{Année 2025-2026}

\begin{document}

\maketitle
\tableofcontents
\newpage

%=============================================================================
\section{Introduction}

Le Set Packing Problem (SPP) est un problème d'optimisation combinatoire NP-difficile.

\subsection{Définition}

Maximiser :
\begin{equation}
    \max_{x \in \{0,1\}^n} \sum_{j=1}^{n} c_j x_j
\end{equation}

Sous contraintes :
\begin{equation}
    \sum_{j=1}^{n} a_{ij} x_j \leq 1, \quad \forall i \in \{1,\ldots,m\}
\end{equation}

%=============================================================================
\section{EI1 : Heuristiques Constructives}

\subsection{Heuristique gloutonne}

Construction itérative basée sur l'utilité :
\begin{equation}
    u_j = \frac{c_j}{\sum_{i=1}^{m} a_{ij}}
\end{equation}

\subsection{Recherche locale}

Descente profonde avec voisinages :
\begin{itemize}
    \item \textbf{1-1} : retirer 1, ajouter 1
    \item \textbf{2-1} : retirer 2, ajouter 1
\end{itemize}

\subsection{Résultats EI1}

% Tableau généré automatiquement
\begin{table}[htbp]
\centering
\caption{Résultats EI1 : Glouton + Recherche Locale}
\label{tab:ei1}
\footnotesize
\begin{tabular}{|l|rr|rrr|rrr|rrrr|}
\hline
\textbf{Instance} & \textbf{m} & \textbf{n} & \textbf{Z*} & \textbf{T(s)} & \textbf{Statut} & \textbf{Z} & \textbf{T(s)} & \textbf{Gap\%} & \textbf{Z} & \textbf{T(s)} & \textbf{Gap\%} & \textbf{Speedup} \\
\hline
didactic.dat & 7 & 9 & --- & --- & N/A & 30 & 0.0 & --- & 30 & 0.0 & --- & --- \\
pb\_100rnd0100.dat & 500 & 100 & 372 & 0.64 & OPTIMAL & 342 & 0.0022 & 8.06 & 348 & 0.0114 & 6.45 & 56.7 \\
pb\_100rnd0300.dat & 500 & 100 & --- & --- & N/A & 193 & 0.0015 & --- & 194 & 0.0035 & --- & --- \\
pb\_200rnd0100.dat & 1000 & 200 & --- & --- & N/A & 351 & 0.0083 & --- & 359 & 0.0251 & --- & --- \\
pb\_200rnd0300.dat & 1000 & 200 & --- & --- & N/A & 682 & 0.0192 & --- & 688 & 0.1595 & --- & --- \\
pb\_500rnd1500.dat & 1500 & 500 & --- & --- & N/A & 1059 & 0.1007 & --- & 1086 & 0.2073 & --- & --- \\
pb\_500rnd0100.dat & 2500 & 500 & --- & --- & N/A & 285 & 0.0375 & --- & 285 & 0.0466 & --- & --- \\
pb\_500rnd0300.dat & 2500 & 500 & --- & --- & N/A & 674 & 0.1939 & --- & 674 & 0.2291 & --- & --- \\
pb\_1000rnd0100.dat & 5000 & 1000 & --- & --- & N/A & 49 & 0.0847 & --- & 49 & 0.0865 & --- & --- \\
pb\_2000rnd0100.dat & 10000 & 2000 & --- & --- & N/A & 37 & 0.5496 & --- & 37 & 0.5572 & --- & --- \\
\hline
\end{tabular}
\end{table}


\begin{figure}[htbp]
    \centering
    \includegraphics[width=0.95\textwidth]{res/graphiques_EI1.pdf}
    \caption{Résultats EI1}
    \label{fig:ei1}
\end{figure}

\subsection{Analyse EI1}

\paragraph{Qualité :} 
Le gap moyen après recherche locale est de \textbf{XX\%}.

\paragraph{Temps :}
Speedup moyen de \textbf{XXx} par rapport à GLPK.

%=============================================================================
\section{EI2 : GRASP et Reactive GRASP}

\subsection{GRASP Standard}

Phase construction avec RCL :
\begin{equation}
    \text{RCL} = \{j : u_j \geq u_{\min} + \alpha (u_{\max} - u_{\min})\}
\end{equation}

\subsection{Reactive GRASP}

Adaptation automatique de $\alpha$ par qualité :
\begin{equation}
    p_k = \frac{q_k}{\sum_{i} q_i}, \quad q_k = \frac{\bar{z}_k - z_{\text{worst}}}{z_{\text{best}} - z_{\text{worst}}}
\end{equation}

\subsection{Résultats EI2}

% Tableau généré automatiquement
\begin{table}[htbp]
\centering
\caption{Résultats EI2 : GRASP vs Reactive GRASP (5 runs)}
\label{tab:ei2}
\small
\begin{tabular}{|l|rr|rrrr|rrrr|}
\hline
\textbf{Instance} & \textbf{m} & \textbf{n} & \multicolumn{4}{c|}{\textbf{GRASP}} & \multicolumn{4}{c|}{\textbf{Reactive GRASP}} \\
& & & Z_{min} & Z_{max} & Z_{moy} & T(s) & Z_{min} & Z_{max} & Z_{moy} & T(s) \\
\hline
didactic.dat & 7 & 9 & 30 & 30 & 30.0 & 0.0 & 30 & 30 & 30.0 & 0.0 \\
pb\_100rnd0100.dat & 500 & 100 & 366 & 368 & 367.6 & 0.39 & 365 & 370 & 367.2 & 0.38 \\
pb\_100rnd0300.dat & 500 & 100 & 195 & 203 & 196.6 & 0.24 & 195 & 203 & 196.6 & 0.23 \\
pb\_200rnd0100.dat & 1000 & 200 & 399 & 403 & 399.8 & 1.09 & 387 & 403 & 394.6 & 1.09 \\
pb\_200rnd0300.dat & 1000 & 200 & 693 & 703 & 698.6 & 2.12 & 689 & 699 & 694.8 & 2.09 \\
pb\_500rnd0100.dat & 2500 & 500 & 289 & 300 & 294.8 & 4.03 & 285 & 305 & 293.4 & 3.91 \\
pb\_500rnd0300.dat & 2500 & 500 & 681 & 714 & 696.4 & 9.23 & 711 & 739 & 723.2 & 9.16 \\
pb\_500rnd1500.dat & 1500 & 500 & 1077 & 1111 & 1088.8 & 10.5 & 1097 & 1113 & 1103.6 & 10.31 \\
pb\_1000rnd0100.dat & 5000 & 1000 & 54 & 67 & 60.6 & 7.58 & 54 & 59 & 57.0 & 7.86 \\
pb\_2000rnd0100.dat & 10000 & 2000 & 39 & 40 & 39.6 & 41.37 & 37 & 40 & 38.8 & 43.12 \\
\hline
\end{tabular}
\end{table}


\begin{figure}[htbp]
    \centering
    \includegraphics[width=0.95\textwidth]{res/graphiques_EI2.pdf}
    \caption{Comparaison GRASP vs Reactive GRASP}
    \label{fig:ei2}
\end{figure}

\subsection{Analyse EI2}

\paragraph{GRASP vs Reactive :}
Reactive GRASP obtient de meilleurs résultats dans \textbf{XX\%} des cas.

\paragraph{Distribution $\alpha$ :}
Les valeurs moyennes ($\alpha = 0.5$-$0.75$) sont privilégiées.

%=============================================================================
\section{EI3 : ACO vs AG vs GRASP}

% À compléter avec ACO ou autre

%=============================================================================
\section{Conclusion}

Les heuristiques développées offrent un excellent compromis qualité/temps.

%=============================================================================
\appendix
\section{Code}

\subsection{Structure}

\begin{verbatim}
src/
├── livrable.jl        # Fichier principal
├── Glouton.jl
├── Exploration.jl
├── Grasp.jl
└── ...

dat/                   # 10 instances
res/                   # Résultats générés
doc/                   # Rapport
\end{verbatim}

\subsection{Utilisation}

\begin{verbatim}
include("src/livrable.jl")

# EI1
experimentationSPP()

# EI2
experimentationSPP_EI2(n_runs=5)
\end{verbatim}

\end{document}